\documentclass[10pt,a4paper]{amsart}
\usepackage[margin=19mm]{geometry}
\usepackage{amsmath,amssymb,mathtools}
\usepackage[scr=rsfs,cal=euler]{mathalfa}
\usepackage{xfrac}
\usepackage[unicode,hidelinks,bookmarks=false]{hyperref}
\newcommand{\ie}{i.e.\ }
\newcommand{\cf}{cf.\ }
\newcommand{\resp}{respectively\ }
\newcommand{\Subsection}{\subsection}
\providecommand{\nobreakdash}{\nobreak\hbox{-}}
\setlength{\emergencystretch}{2em}
\multlinegap0pt
\usepackage{float}
\usepackage{amssymb}
\usepackage{enumerate}
\usepackage{csquotes}
\usepackage{mathtools}
\usepackage{xcolor}
\newtheorem{theorem}{Theorem}
\title{A Study on Cumulative Residual Extropy of Linear Consecutive k-out-of-n:G Systems}
\author{Aman Pandey, Chanchal Kundu}
\date{Source: arXiv:2607.06992v1 (CC BY 4.0)}
\begin{document}
\maketitle

\noindent\emph{Source and licence.} A Study on Cumulative Residual Extropy of Linear Consecutive k-out-of-n:G Systems. Version arXiv:2607.06992v1 (CC BY 4.0), distributed by arXiv under the Creative Commons Attribution 4.0 International licence, http://creativecommons.org/licenses/by/4.0/, as recorded in the arXiv licence metadata for this exact version and shown on the abstract page https://arxiv.org/abs/2607.06992v1. Full licence notice: \texttt{licenses/arxiv-2607.06992-CC-BY-4.0.txt}.\par\smallskip

\noindent\textit{Editorial scope.} Counted unit: the article's theorem theorem (source0.tex lines 299--306). The article's complete original proof of it is delivered with it (source0.tex lines 307--433, one \texttt{proof} environment of 127 lines). No mathematics is written, repaired or re-derived: the statement and the proof are the article's own lines, in the article's own notation, and no retained line is substituted or re-flowed.  Retained uncounted as the source-local context the proof needs: lines 262--267.  Nothing was omitted from any retained span, so the package carries no omission record.  No retained line contains a citation, so the wrapper carries no bibliography and no bibliography entry.  No retained line contains a figure environment, an image-inclusion command or drawing code, so no image is delivered and the package declares no asset dependency.  Editorial formatting only, in addition: an AMS \texttt{amsart} preamble that restates the article's own declarations and macros used by the retained lines, namely the environments \texttt{theorem} on separate counters, together with the standard packages those lines need.  The retained environments print in this document as Theorem 1 in order of appearance, whereas the article prints the same items with the numbering of its own full document.  The complete native source of the article as distributed in the arXiv e-print for this exact version is bundled unchanged as \texttt{sources/204675/source0.tex}.
\begin{align}\label{eq2.4}
    X\leq_{lir}Y
    \iff
    \eta_Y^{-1}(w)-\eta_X^{-1}(w)
    \text{ is increasing in } w>0.
\end{align}

\begin{theorem}
Let $X\leq_{lir}Y$. Further, assume that
\begin{align*}
    \Phi(t):=
    \frac{\Phi(\bar G_{k|n:G}(s))}{s},~0<s<1,
\end{align*}
is decreasing in $s$. Then, $J(T_{k|n:G}^X)\leq J(T_{k|n:G}^Y).$
\end{theorem}

\begin{proof}
Recall that
\begin{align*}
    \bar G_{k|n:G}(v)
    =
    (n-k+1)(1-v)^k-(n-k)(1-v)^{k+1},
    \qquad 0<v<1.
\end{align*}
Hence, the SF  can be expressed as
$S^X_{k|n:G}(w)
    =
    \bar G_{k|n:G}(F_X(w)).$ Since $X\leq_{lir}Y$, relation \eqref{eq2.4} yields $\eta_Y^{-1}(w)-\eta_X^{-1}(w)$
is increasing in $w>0$. Differentiating both sides gives
\begin{align*}
    \frac{d}{dw}
    \left[
    \eta_Y^{-1}(w)-\eta_X^{-1}(w)
    \right]
    =
    \frac{1}{F_Y(\eta_Y^{-1}(w))}
    -
    \frac{1}{F_X(\eta_X^{-1}(w))}
    \geq 0,
\end{align*}
which implies
\begin{align}\label{eq:new1}
    F_X(w)
    \geq
    F_Y\!\left(
    \eta_Y^{-1}(\eta_X(w))
    \right),~ w>0.
\end{align}

Now, by the definition of CREx,
\begin{align}
    J(T^X_{k|n:G})
    &=
    -\frac12
    \int_0^\infty
    \left(
    S^X_{k|n:G}(w)
    \right)^2
    dw
    \nonumber\\
    &=-\frac12
    \int_0^\infty
    \Phi\!\left(
    \bar G_{k|n:G}(F_X(w))
    \right)
    dw
    \nonumber\\
    &=-\frac12
    \int_0^\infty
    \frac{
    \Phi\!\left(
    \bar G_{k|n:G}(F_X(w))
    \right)
    }{F_X(w)}
    F_X(w)\,dw.
    \label{eq:new2}
\end{align}
Since $\frac{\Phi(\bar G_{k|n:G}(s))}{s}$
is decreasing in $s$, together with \eqref{eq:new1}, we obtain
\begin{align}
    \frac{
    \Phi(\bar G_{k|n:G}(F_X(w)))
    }{F_X(w)}
    \leq
    \frac{
    \Phi\!\left(
    \bar G_{k|n:G}
    \left(
    F_Y(\eta_Y^{-1}(\eta_X(w)))
    \right)
    \right)
    }{
    F_Y(\eta_Y^{-1}(\eta_X(w)))
    }.
    \label{eq:new3}
\end{align}

Thus, from \eqref{eq:new3} and \eqref{eq:new2}, we have
\begin{align}
    J(T^X_{k|n:G})
    \geq -\frac12
    \int_0^\infty
    \frac{
    \Phi\!\left(
    \bar G_{k|n:G}
    \left(
    F_Y(\eta_Y^{-1}(\eta_X(w)))
    \right)
    \right)
    }{
    F_Y(\eta_Y^{-1}(\eta_X(w)))
    }
    F_X(w)\,dw.
    \label{eq:new4}
\end{align}
Now, perform the transformation $v=\eta_Y^{-1}(\eta_X(w)).$
Then, 
$\eta_Y(v)=\eta_X(w).$
Differentiating both sides yields, 
$dw
    =
    \frac{F_Y(v)}{F_X(w)}\,dv.$

Therefore, \eqref{eq:new4} becomes
\begin{align*}
    J(T^X_{k|n:G})
    &\geq -\frac12
    \int_{\eta_Y^{-1}(\eta_X(0))}^{\infty}
    \Phi\!\left(
    \bar G_{k|n:G}(F_Y(v))
    \right)\,dv
    \\
    &=
    -\frac12
    \int_0^\infty
    \left(
    S^Y_{k|n:G}(v)
    \right)^2
    dv=
    J(T^Y_{k|n:G}).
\end{align*}
This completes the proof.
\end{proof}

\end{document}
